\section{Koeficient determinace}\label{sec:ai-r2}

Samotná hodnota \(\MSE\) závisí na měřítku výstupního souboru. Abychom chybu regresní přímky zasadili do kontextu, porovnáme ji s~nejlepším konstantním modelem
\[
    c(x)=\average{\mathbf{Y}}.
\]
Podle vlastností průměru platí
\[
    \MSE_{\mathbf{X},\mathbf{Y}}(c)
    =
    \frac1n\sum_{i=1}^{n}
    (y_i-\average{\mathbf{Y}})^2
    =
    \Var(\mathbf{Y}).
\]

\begin{definition}\label{def:ai-r2}
    Nechť \(\sigma_{\mathbf{Y}}>0\) a~\(f\) je regresní přímka souborů
    \(\mathbf{X},\mathbf{Y}\). \emph{Koeficient determinace} definujeme jako
    \[
        R^2
        =
        1-
        \frac{
            \MSE_{\mathbf{X},\mathbf{Y}}(f)
        }{
            \MSE_{\mathbf{X},\mathbf{Y}}
            (x\mapsto\average{\mathbf{Y}})
        }
        =
        1-
        \frac{
            \MSE_{\mathbf{X},\mathbf{Y}}(f)
        }{
            \Var(\mathbf{Y})
        }.
    \]
\end{definition}

Pro jednoduchou lineární regresi dosadíme vztah z~
\smartref[tvrzení~][]{\showrefs}{prop:ai-regrese-min-mse}{tvrzení o~minimální MSE}:
\[
    R^2
    =
    1-
    \frac{
        \sigma_{\mathbf{Y}}^2
        (1-\rho_{\mathbf{X}\mathbf{Y}}^2)
    }{
        \sigma_{\mathbf{Y}}^2
    }
    =
    \rho_{\mathbf{X}\mathbf{Y}}^2.
\]
Proto \(0\leqslant R^2\leqslant1\). Hodnota blízká jedné znamená, že regresní přímka odstranila velkou část čtvercové chyby konstantního modelu. Hodnota blízká nule znamená, že lineární závislost přinesla jen malé zlepšení.

\begin{figure}[ht]
    \centering
    \begin{tikzpicture}
    \begin{axis}[
        width=7.2cm,height=5.5cm,
        xmin=0,xmax=6,ymin=0,ymax=6,
        axis lines=left,xtick=\empty,ytick=\empty,
        title={vysoké \(R^2\)},
        title style={font=\normalsize},
        every node/.style={font=\normalsize}
    ]
        \addplot[red!75!black,very thick,domain=0.4:5.6]{0.8*x+0.5};
        \addplot[only marks,mark=*,blue!75!black]
        coordinates{(0.7,1.1)(1.4,1.5)(2.1,2.3)(2.8,2.6)(3.5,3.5)(4.2,3.8)(5,4.7)};
    \end{axis}
    \begin{axis}[
        xshift=8.2cm,
        width=7.2cm,height=5.5cm,
        xmin=0,xmax=6,ymin=0,ymax=6,
        axis lines=left,xtick=\empty,ytick=\empty,
        title={nižší \(R^2\)},
        title style={font=\normalsize},
        every node/.style={font=\normalsize}
    ]
        \addplot[red!75!black,very thick,domain=0.4:5.6]{0.55*x+1.2};
        \addplot[only marks,mark=*,orange!85!black]
        coordinates{(0.7,0.7)(1.3,2.8)(2.1,1.3)(2.8,3.5)(3.4,2.1)(4.2,4.6)(5,3.2)};
    \end{axis}
\end{tikzpicture}

    \caption{Stejný trend může být daty zachycen s~různou přesností.}
    \label{fig:ai-r2}
\end{figure}

\begin{example}
    Pro
    \[
        \mathbf{X}=(1,2,3,4,5),
        \qquad
        \mathbf{Y}=(2,3,5,4,6)
    \]
    dostaneme
    \[
        f(x)=0{,}9x+1{,}3,
        \qquad
        \MSE_{\mathbf{X},\mathbf{Y}}(f)=0{,}38.
    \]
    Konstantní model má chybu
    \[
        \Var(\mathbf{Y})=2,
    \]
    a~tedy
    \[
        R^2=1-\frac{0{,}38}{2}=0{,}81.
    \]
    Regresní přímka odstranila \(81\,\%\) střední kvadratické chyby nejlepšího konstantního modelu.
\end{example}

\warningbox{Vysoké \(R^2\) nedokazuje příčinný vztah ani správnost modelu mimo rozsah dat. V~obecnějších regresních úlohách se navíc vztah \(R^2=\rho_{\mathbf{X}\mathbf{Y}}^2\) nemusí použít v~této jednoduché podobě.}
